\documentclass[aspectratio=169,10pt]{beamer}\usetheme{metropolis}
\title{Q04: RL-Based Reasoning vs. Ordinary Generation}\subtitle{How the optimization signal changes}\author{DLS 2026}\date{October 2026}
\begin{document}\maketitle
\begin{frame}{The question}\begin{block}{Q04}How does RL-based reasoning differ from ordinary LLM generation?\end{block}\begin{itemize}\item Generation: sample from a fixed policy.\item Training: change the policy using an objective.\end{itemize}\end{frame}
\begin{frame}{Pre-training and inference}\begin{itemize}\item Pre-training: next-token cross-entropy on a corpus.\item Inference: $\theta$ is fixed; prompt and decoding influence samples.\item Longer reasoning text can be elicited without a parameter update.\end{itemize}\end{frame}
\begin{frame}{Post-training adds a target}\begin{columns}\column{.32\textwidth}\textbf{SFT}\\imitate selected responses\column{.32\textwidth}\textbf{Preference optimization}\\favor chosen over rejected responses\column{.32\textwidth}\textbf{RL}\\maximize expected reward\end{columns}\vfill These are mechanisms, not guarantees of reasoning.\end{frame}
\begin{frame}{Idealized RL objective}\[\max_\theta\;\mathbb E_{y\sim\pi_\theta}[r(x,y)]-\beta D_{\mathrm{KL}}(\pi_\theta\|\pi_{\mathrm{ref}})\]\begin{itemize}\item Reward says which outcome is preferred.\item KL term constrains policy drift.\item The practical result depends on the reward, data, and optimization.\end{itemize}\end{frame}
\begin{frame}{Why verifiable rewards matter}\begin{itemize}\item Tests, proof checkers, and numeric oracles can score outcomes.\item Process supervision can score intermediate steps.\item A verifier is only as good as its specification and coverage.\end{itemize}\alert{Reward what you can verify; audit what you failed to verify.}\end{frame}
\begin{frame}{Do not equate reward with truth}\begin{itemize}\item Reward hacking and shortcuts remain possible.\item A correct final answer need not imply a faithful rationale.\item Held-out, adversarial, and cost-controlled evaluation is required.\end{itemize}\end{frame}
\begin{frame}{Toy experiment}\begin{itemize}\item Sample one of several candidate answers.\item Give reward 1 only to the verifier-approved answer.\item Update probabilities; inspect how mass shifts.\end{itemize}\vfill A bandit demonstration, not RLHF or an LLM.\end{frame}
\begin{frame}{Takeaway}\begin{block}{Core difference}Ordinary generation samples a learned distribution. RL-based post-training changes that distribution using a reward/preference objective. The reward defines the behavior being optimized---and its blind spots.\end{block}\vfill\small References: Ouyang et al. (2022); Rafailov et al. (2023); Lightman et al. (2023).\end{frame}
\end{document}
